\documentclass{article}
\usepackage{amsmath,amssymb,amsthm}
\usepackage{algorithm,algorithmic}
\usepackage{bm}
\newtheorem{theorem}{Theorem}
\newtheorem{lemma}{Lemma}
\newcommand{\E}{\mathbb{E}}
\newcommand{\norm}[1]{\left\|#1\right\|}
\newcommand{\ip}[2]{\left\langle #1,#2\right\rangle}
\begin{document}

\section*{Quantized Stochastic Gradient Descent (QSGD)}

\section*{Mathematical Formulation}
Let $f:\mathbb{R}^d\rightarrow\mathbb{R}$ be a possibly non‑convex, $L$‑smooth function.  
At iteration $t$ we draw an i.i.d. data sample $\xi_t$ and form a stochastic gradient
\[
g_t \;:=\; \nabla f(w_t; \xi_t).
\]
A (possibly biased) quantization operator $Q:\mathbb{R}^d\rightarrow\mathbb{R}^d$ satisfies  

\[
\boxed{\E[\,Q(g_t)\mid w_t\,]=g_t},\qquad 
\boxed{\E\big[\|Q(g_t)-g_t\|^2\mid w_t\big]\le \sigma_q^2},
\tag{1}
\]
i.e. it is unbiased and has bounded second moment.  

The update rule of QSGD is
\[
\boxed{w_{t+1}=w_t-\eta_t\,Q(g_t)},
\tag{2}
\]
where $\{\eta_t\}_{t\ge0}$ is a stepsize sequence (constant or diminishing).

\section*{Pseudocode}
\begin{algorithm}[H]
\caption{Quantized Stochastic Gradient Descent}
\begin{algorithmic}[1]
\STATE \textbf{Input:} initial point $w_0\in\mathbb{R}^d$, stepsizes $\{\eta_t\}$, quantizer $Q(\cdot)$
\FOR{$t=0,1,\dots,T-1$}
\STATE Sample a minibatch (or a single data point) $\xi_t$
\STATE Compute stochastic gradient $g_t=\nabla f(w_t; \xi_t)$
\STATE Quantize: $\tilde g_t = Q(g_t)$
\STATE Update: $w_{t+1}=w_t-\eta_t \tilde g_t$
\ENDFOR
\STATE \textbf{Return} $w_T$ (or the averaged iterate $\bar w_T = \frac1T\sum_{t=0}^{T-1} w_t$)
\end{algorithmic}
\end{algorithm}

\section*{Dry Run Example}
Consider the scalar quadratic $f(w)=(w-3)^2$, thus $\nabla f(w)=2(w-3)$.  
Take a deterministic “quantizer’’ that returns the exact gradient with probability $1$ (so $Q(g)=g$) – this satisfies (1) with $\sigma_q=0$.  

\begin{center}
\begin{tabular}{c|c|c|c|c}
$t$ & $w_t$ & $g_t= \nabla f(w_t)$ & $Q(g_t)$ & $w_{t+1}=w_t-\eta g_t$ \\ \hline
0 & $0$ & $-6$ & $-6$ & $0-0.1(-6)=0.6$ \\
1 & $0.6$ & $-4.8$ & $-4.8$ & $0.6-0.1(-4.8)=1.08$ \\
2 & $1.08$ & $-3.84$ & $-3.84$ & $1.08-0.1(-3.84)=1.464$ \\
3 & $1.464$ & $-3.072$ & $-3.072$ & $1.464-0.1(-3.072)=1.7712$
\end{tabular}
\end{center}
Corresponding objective values:
$f(0)=9,\; f(0.6)=5.76,\; f(1.08)=3.6864,\; f(1.7712)=1.5625$,
showing descent toward the optimum $w^\star=3$.

\newpage
\section*{Assumptions}
\begin{enumerate}
\item \textbf{Smoothness.} $f$ is $L$‑smooth:
\[
\forall x,y\in\mathbb{R}^d,\qquad 
f(y)\le f(x)+\ip{\nabla f(x)}{y-x}+\frac{L}{2}\norm{y-x}^2.
\tag{A1}
\]
\item \textbf{Unbiased stochastic gradient.}
\[
\E\big[g_t\mid w_t\big]=\nabla f(w_t),\qquad
\E\big[\|g_t-\nabla f(w_t)\|^2\mid w_t\big]\le\sigma_g^2.
\tag{A2}
\]
\item \textbf{Unbiased quantization with bounded variance.} Equation (1) holds with a constant $\sigma_q\ge0$.
\item \textbf{Stepsize.} Either a constant $\eta_t=\eta\le \frac{1}{L}$ or a diminishing schedule
\[
\eta_t=\frac{c}{\sqrt{t+1}},\qquad c>0.
\tag{A3}
\]
\item \textbf{Bounded iterates (optional).} $\norm{w_t}\le R$ for all $t$, which holds automatically for many practical choices of $\eta$ because of the smoothness bound.
\end{enumerate}

\section*{Main Convergence Theorem}
\begin{theorem}[Convergence to Stationarity]\label{thm:main}
Under Assumptions (A1)–(A3) and using the constant stepsize $\eta\le\frac{1}{L}$,
the iterates generated by \eqref{2} satisfy
\[
\frac{1}{T}\sum_{t=0}^{T-1}\E\!\big[\|\nabla f(w_t)\|^2\big]
\;\le\;
\frac{2\big(f(w_0)-f^\star\big)}{\eta T}
\;+\;
L\big(\sigma_g^2+\sigma_q^2\big),
\tag{3}
\]
where $f^\star=\inf_{w}f(w)$. Consequently,
\[
\min_{0\le t<T}\E\!\big[\|\nabla f(w_t)\|^2\big]=O\!\Big(\frac{1}{T}\Big)+O\!\big(L(\sigma_g^2+\sigma_q^2)\big).
\]
\end{theorem}

\section*{Theoretical Properties}
\begin{itemize}
\item \textbf{Stability.} The unbiasedness of $Q$ guarantees that the expected descent direction coincides with the true stochastic gradient, thus the algorithm inherits the stability of standard SGD.
\item \textbf{Hyper‑parameter sensitivity.} The bound \eqref{3} shows a linear dependence on the stepsize $\eta$ in the first term and a quadratic dependence on $L$ in the variance term. Choosing $\eta\approx \frac{1}{L}$ balances both contributions.
\item \textbf{Comparison to baseline SGD.} If $Q$ is the identity operator ($\sigma_q=0$) the bound reduces to the classic SGD result. Quantization adds an additive variance term $L\sigma_q^2$.
\item \textbf{Effect of diminishing stepsizes.} With $\eta_t=c/\sqrt{t+1}$ the first term in \eqref{3} decays as $O(1/\sqrt{T})$ while the variance term remains $O(L(\sigma_g^2+\sigma_q^2))$, yielding the standard $O(1/\sqrt{T})$ rate for non‑convex stochastic optimization.
\end{itemize}

\section*{Discussion}
The theorem demonstrates that quantization does not alter the order of convergence; it merely inflates the steady‑state variance by $L\sigma_q^2$. In practice, modern stochastic quantizers are designed so that $\sigma_q^2$ scales as $O(\|g_t\|^2 /s^2)$ where $s$ is the number of quantization levels, making the extra term negligible for sufficiently fine quantization.

\newpage
\section*{Preliminaries (Definitions and Lemmas)}
\begin{lemma}[Smoothness inequality]\label{lem:smooth}
If $f$ is $L$‑smooth, then for any $x,y$,
\[
f(y)\le f(x)+\ip{\nabla f(x)}{y-x}+\frac{L}{2}\norm{y-x}^2.
\]
\end{lemma}
\begin{lemma}[Unbiasedness of composite noise]\label{lem:noise}
Define the total noise $\xi_t:=Q(g_t)-\nabla f(w_t)$. Under (A2) and (A1) we have
\[
\E[\xi_t\mid w_t]=0,\qquad 
\E\big[\|\xi_t\|^2\mid w_t\big]\le\sigma_g^2+\sigma_q^2.
\]
\end{lemma}
\begin{proof}[Proof of Lemma \ref{lem:noise}]
\[
\E[\xi_t\mid w_t]
= \E[Q(g_t)-\nabla f(w_t)\mid w_t]
= \E[Q(g_t)\mid w_t]-\nabla f(w_t)
= \E[\E[Q(g_t)\mid g_t,w_t]\mid w_t]-\nabla f(w_t)
= \E[g_t\mid w_t]-\nabla f(w_t)=0.
\]
For the second moment,
\[
\begin{aligned}
\E\!\big[\|\xi_t\|^2\mid w_t\big]
&=\E\!\big[\|Q(g_t)-g_t+g_t-\nabla f(w_t)\|^2\mid w_t\big]\\
&\overset{[Step 1]}{\le}
2\E\!\big[\|Q(g_t)-g_t\|^2\mid w_t\big]
+2\E\!\big[\|g_t-\nabla f(w_t)\|^2\mid w_t\big]\\
&\le 2\sigma_q^2+2\sigma_g^2
= \sigma_g^2+\sigma_q^2\quad\text{(renaming the constant).}
\end{aligned}
\]
\end{proof}

\section*{Main Proof (Step‑by‑Step Derivation)}
We prove Theorem \ref{thm:main} for a constant stepsize $\eta\le\frac{1}{L}$.
\subsection*{Step 1 – Smoothness descent}
Using Lemma \ref{lem:smooth} with $x=w_t$, $y=w_{t+1}=w_t-\eta Q(g_t)$,
\[
f(w_{t+1})\le f(w_t)-\eta\ip{\nabla f(w_t)}{Q(g_t)}+\frac{L\eta^2}{2}\norm{Q(g_t)}^2.
\tag{4}
\]

\subsection*{Step 2 – Expand the inner product}
Write $Q(g_t)=\nabla f(w_t)+\xi_t$ where $\xi_t$ is the total noise from Lemma \ref{lem:noise}. Then
\[
\ip{\nabla f(w_t)}{Q(g_t)}
=\norm{\nabla f(w_t)}^2+\ip{\nabla f(w_t)}{\xi_t}.
\tag{5}
\]

\subsection*{Step 3 – Expand the squared norm}
\[
\norm{Q(g_t)}^2
=\norm{\nabla f(w_t)+\xi_t}^2
=\norm{\nabla f(w_t)}^2+2\ip{\nabla f(w_t)}{\xi_t}+\norm{\xi_t}^2.
\tag{6}
\]

\subsection*{Step 4 – Plug (5)–(6) into (4)}
\[
\begin{aligned}
f(w_{t+1})
&\le f(w_t)-\eta\big(\norm{\nabla f(w_t)}^2+\ip{\nabla f(w_t)}{\xi_t}\big)\\
&\qquad+\frac{L\eta^2}{2}\big(\norm{\nabla f(w_t)}^2+2\ip{\nabla f(w_t)}{\xi_t}+\norm{\xi_t}^2\big).
\end{aligned}
\tag{7}
\]

\subsection*{Step 5 – Collect like terms}
\[
\begin{aligned}
f(w_{t+1})\le
&f(w_t)
-\Big(\eta-\frac{L\eta^2}{2}\Big)\norm{\nabla f(w_t)}^2\\
&-\Big(\eta-\!L\eta^2\Big)\ip{\nabla f(w_t)}{\xi_t}
+\frac{L\eta^2}{2}\norm{\xi_t}^2.
\end{aligned}
\tag{8}
\]

\subsection*{Step 6 – Take conditional expectation $\E[\cdot\mid w_t]$}
Using $\E[\xi_t\mid w_t]=0$ (Lemma \ref{lem:noise}),
\[
\E\!\big[f(w_{t+1})\mid w_t\big]
\le f(w_t)-\Big(\eta-\frac{L\eta^2}{2}\Big)\norm{\nabla f(w_t)}^2
+\frac{L\eta^2}{2}\E\!\big[\norm{\xi_t}^2\mid w_t\big].
\tag{9}
\]

\subsection*{Step 7 – Upper‑bound the noise variance}
From Lemma \ref{lem:noise},
\[
\E\!\big[\norm{\xi_t}^2\mid w_t\big]\le\sigma_g^2+\sigma_q^2.
\tag{10}
\]

\subsection*{Step 8 – Simplify the coefficient}
Because $\eta\le\frac{1}{L}$,
\[
\eta-\frac{L\eta^2}{2}\;\ge\;\frac{\eta}{2}.
\tag{11}
\]

\subsection*{Step 9 – Combine (9)–(11)}
\[
\E\!\big[f(w_{t+1})\mid w_t\big]
\le f(w_t)-\frac{\eta}{2}\norm{\nabla f(w_t)}^2
+\frac{L\eta^2}{2}\big(\sigma_g^2+\sigma_q^2\big).
\tag{12}
\]

\subsection*{Step 10 – Take full expectation}
Denote $\Delta_t:=\E[f(w_t)]-f^\star$, where $f^\star:=\inf_w f(w)$.  
Taking expectation of (12) yields
\[
\Delta_{t+1}\le\Delta_t-\frac{\eta}{2}\E\!\big[\norm{\nabla f(w_t)}^2\big]
+\frac{L\eta^2}{2}\big(\sigma_g^2+\sigma_q^2\big).
\tag{13}
\]

\subsection*{Step 11 – Sum from $t=0$ to $T-1$}
\[
\sum_{t=0}^{T-1}\E\!\big[\norm{\nabla f(w_t)}^2\big]
\le\frac{2}{\eta}\big(\Delta_0-\Delta_T\big)
+L\eta\,T\big(\sigma_g^2+\sigma_q^2\big).
\tag{14}
\]

Since $\Delta_T\ge0$, we drop it and divide by $T$:
\[
\frac{1}{T}\sum_{t=0}^{T-1}\E\!\big[\norm{\nabla f(w_t)}^2\big]
\le\frac{2\Delta_0}{\eta T}+L\eta\big(\sigma_g^2+\sigma_q^2\big).
\tag{15}
\]

\subsection*{Step 12 – Choose $\eta$}
Setting $\eta=\frac{1}{L}$ (the largest admissible constant stepsize) gives the bound
\[
\frac{1}{T}\sum_{t=0}^{T-1}\E\!\big[\norm{\nabla f(w_t)}^2\big]
\le\frac{2L\big(f(w_0)-f^\star\big)}{T}+L\big(\sigma_g^2+\sigma_q^2\big),
\]
which is exactly inequality \eqref{3}.  ∎

\section*{Rate Derivation (With Constants)}
From \eqref{15} we see two regimes:

\begin{enumerate}
\item \textbf{Transient regime} ($T$ small): the term $\frac{2\Delta_0}{\eta T}$ dominates and we obtain an $O(1/T)$ decay of the average squared gradient norm.
\item \textbf{Steady‑state regime} ($T\to\infty$): the variance term $L\eta(\sigma_g^2+\sigma_q^2)$ persists.  Choosing a diminishing stepsize $\eta_t=c/\sqrt{t+1}$ makes this term decay as $O(1/\sqrt{T})$, reproducing the classic $O(1/\sqrt{T})$ rate for non‑convex stochastic optimization.
\end{enumerate}

\section*{Special Cases}
\begin{itemize}
\item \textbf{No quantization noise} ($\sigma_q=0$): The bound reduces to the standard SGD result.
\item \textbf{Deterministic gradients} ($\sigma_g=0$): Only quantization contributes to the variance term.
\item \textbf{High‑precision quantizer} where $\sigma_q^2\le\alpha\,\E\!\big[\|g_t\|^2\big]$ with $\alpha\ll1$: The steady‑state term becomes $L\eta\alpha\,\E\!\big[\|g_t\|^2\big]$, which can be absorbed into the transient term, yielding near‑exact SGD behaviour.
\end{itemize}

\end{document}